\subsection{Entitäten}

Sei mit $E$ ein Typ gegeben, der eine Wertemenge\footnote{
    Beispielswiese \texttt{uint32\_t} in \texttt{C++}.
} zur Identifikation zusammengehöriger Komponentenbündel ermöglicht.\\

\noindent
Wenn der Term $e$ einen bestimmten Wert des Typs $E$ repräsentiert, dann ist
\[
    \mathcal{P} \coloneqq E \times \bigcup_{i \in I} K_i = \{(e, k_i) \mid e \in E, i \in I, k_i \in K_i\}
\]
die Menge der möglichen Entität-Komponenten-Paare.\\

\bigskip
Eine bestimmte, eindeutige Ausprägung einer Entität vom Typ $E$, die wir mit $\mathcal{E}_m, m \in \mathbb{N}$ notieren, und die damit verbundene mögliche Menge von Entität-Komponenten-Paaren kann mithilfe folgender partiellen Abbildung hergeleitet werden:
\[
    \mathcal{E}_m :  I \rightharpoonup \bigcup_{i \in I} K_i, \quad i \mapsto k_i \in K_i, m \in \mathbb{N}
\]
Hierbei gilt mit der Definition der partiellen Abbildung:
\[
    \mathcal{E}_m(i) \neq \bot \Leftrightarrow \exists{!} k_i \in K_i : \mathcal{E}_m(i) = k_i
\]

\noindent
Hiermit lässt sich die \textit{assoziierte Komponentenmenge} definieren:

\[
    \mathcal{E}^K_m \coloneqq \{k_i \mid k_i = \mathcal{E}_m(i), k_i \in K_i, i \in I \}
\]

\noindent
\textbf{Beispiel}:

\[
    \begin{alignat}{3}
        \mathcal{E}^K_{1} &= \{\mathcal{E}_{1}(i_1), \mathcal{E}_{1}(i_2), \ldots, \mathcal{E}_{1}(i_n)\}\\
        &= \{k_{i_1}, k_{i_2}, \ldots, k_{i_n}\}
    \end{alignat}
\]

\vspace{5mm}
\noindent
Mit der Definition von $\mathcal{E}_m$ gilt für verschiedene $m, n$ trivialerweise, dass zwei verschiedene Entitäten dieselbe Menge von Termen besitzen können -- damit lässt sich die Eindeutigkeit einer Entität nicht über ihre entsprechende Komponentenmenge feststellen.\\

Aus diesem Grund wird bei Redmond et al. diese Darstellung nicht verwendet.
Ganz im Sinne der Datenorientierung ist die Zuordnung der Entitäten zu Komponenten über Zustände definiert, die eine Menge von Abbildungen der Komponenten-Typen auf eine Menge von partiellen Abbildungen $e_m \rightarrow k_i$ enthält (siehe Abschnitt~\ref{sec:zustand}).
Dies erleichtert auch die formale Herleitung der Identität von Entitäten: Somit exitiert in dem späteren Modell immer nur genau eine eindeutig identifizierbare Entität-Komponentent-Konstellation.
\bigskip

\noindent
Redmond et al. binden deshalb die Eindeutigkeit einer Entität an den Term $e \in E$:

\[
    E \coloneqq \{e_1, e_2, \ldots\}
\]

\noindent
so dass alle Tupel in $\mathcal{P}$, die den Term $e_m$ enthalten, über die Familie

\[
    \mathcal{P}_{e_m} \coloneqq \{(e_m, k_i) \mid e_m \in E, \mathcal{E}_m(i) = k_i \in K_i\}, m \in \mathbb{N}
\]

\noindent
definiert werden kann.

\bigskip
\begin{theorem}\label{th:entity_id1}
    Gegeben sei $\mathcal{P}_{e_m}$ und $\mathcal{P}_{e_n}$, $e_n, e_m \in E$.
    Dann gilt:
    \[
        e_m \neq e_n \Rightarrow \mathcal{P}_{e_m} \neq \mathcal{P}_{e_n}
    \]
\end{theorem}
\bigskip

\begin{proof}
    Seien beliebige $x = (e_m, k_{i}) \in P_{e_m}, y = (e_n, k_j) \in P_{e_n}, k_i \in K_i, j \in K_j, i, j \in I$ gegeben.\\

    \noindent
    Mit der Voraussetzung $e_m \neq e_n$ gilt
    \[
        e_m \neq e_n \land (k_i = k_j \lor k_i \neq k_j)
    \]

    \begin{itemize}
        \item \textbf{Fall 1} - $e_m \neq e_n \land k_i = k_j$: Es ist $x(1) \neq y(1) \land x(2) = y(2)$ und damit $x \neq y$.
        \item \textbf{Fall 2} - $e_m \neq e_n \land k_i \neq k_j$: Es ist $x(1) \neq y(1) \land x(2) \neq y(2)$ und damit $x \neq y$.
    \end{itemize}

    \noindent
    In beiden Fällen ist  $\mathcal{P}_{e_n} \neq \mathcal{P}_{e_m}$ erfüllt.
\end{proof}
\bigskip

\noindent
Theorem~\ref{th:entity_id2} folgt direkt als Kontraposition zu Theorem~\ref{th:entity_id1}:\footnote{
    $A \Rightarrow B \equiv \neg B \Rightarrow \neg A$.
    Wir geben den Beweis trotzdem formal an, um die Bindung der Komponenten-Terme an den Entitätsterm zu verdeutlichen.
}

\begin{theorem}\label{th:entity_id2}
    Gegeben sei $\mathcal{P}_{e_m}$ und $\mathcal{P}_{e_n}$, $e_n, e_m \in E$.
    Dann gilt:
    \[
        \mathcal{P}_{e_m} = \mathcal{P}_{e_n} \Rightarrow e_m = e_n
    \]
\end{theorem}
\bigskip

\begin{proof}
    Seien beliebige $x = (e_m, k_{i}) \in P_{e_m}, y = (e_n, k_j) \in P_{e_n}, k_i \in K_i, j \in K_j, i, j \in I$ gegeben.\\

    \noindent
    Mit der Voraussetzung $\mathcal{P}_{e_m} = \mathcal{P}_{e_n}$ gilt
    \[
        \forall x \in \mathcal{P}_{e_m} : y \in \mathcal{P}_{e_n}
    \]

    \noindent
    Damit folgt direkt
    \[
        \begin{alignat}{3}
            x = y &\Leftrightarrow (e_m, k_i) = (e_n, k_j) \\
            &\Leftrightarrow e_m = e_n \land k_i = k_j
        \end{alignat}
    \]
\end{proof}
\bigskip

\noindent
Wir haben damit eine triviale, aber wichtige Eigenschaft gezeigt: Die Identität einer Entität-Komponenten-Konstellation kann erst mit $e \in E$ definiert werden.
Dies wird später relevant, wenn es darum geht, den nebenläufigen Zugriff auf Entitäten und ihre Komponenten zu formalisieren.\\

Eine weitere wichtige Abgrenzung nehmen Redmond et al. bzgl. der assoziierten Komponententypen vor:

\[
    \forall (e_m, k_i), (e_m, k_j) \in \mathcal{P}_{e_m} : k_i \neq k_j \Rightarrow i \neq j
\]

\noindent
Eine Entität ist also mit höchstens einem Term $k_i$ pro Komponententyp $K_i$ assoziiert.
Davon unberührt bleiben die Werte der Terme: So kann auch weiterhin $k_i = k_j, i \neq j$ gelten.
